\begin{textsample}{NEG Dim 2 – 2019\_c5 – Score -2.00 – 2019\_c5/enade\_2019\_med\_vet\_c996fe1563.txt}  \label{ex:f2_neg_099}
CAPÍTULO 5 DISTRIBUIÇÃO DOS CONCEITOS 5.1 PANORAMA NACIONAL DA DISTRIBUIÇÃO DOS CONCEITOS O Conceito Enade é um indicador de qualidade que avalia os cursos por intermédio dos desempenhos dos estudantes a partir dos resultados do Enade.
Os valores possíveis vão de 1 ( pior situação ) a 5 ( melhor situação ).
Os cursos com apenas um ou sem nenhum concluinte participante não obtêm o Conceito Enade, ficando Sem Conceito ( SC ).
Na Tabela 5.1, são apresentadas a quantidade e a distribuição de cursos de Medicina Veterinária participantes do Enade/2019, por faixa de conceito e Grande Região. \textbf{Enfatiza} -se, mais uma vez, que a diferença entre os cursos tabulados neste capítulo e no capítulo 2 corresponde aos cursos que não foram avaliados, ou seja, àqueles, em princípio, sem alunos concluintes que tivessem sido inscritos para a prova.
Observando -se os dados da Tabela 5.1, nota -se que, dos 215 cursos participantes, 77 ( 35,8 \% ) classificaram -se com conceito 3, o valor modal.
Esse foi também o conceito modal nas regiões : Nordeste ( 51,6 \% ), Sudeste ( 34,4 \% ) e Sul ( 37,3 \% ).
O conceito 2 foi o modal na região Norte ( 35,3 \% ) e os conceitos 1 ou 4 foram os modais na região Centro-Oeste ( 30,4 \%, para cada ).
O conceito 2 foi o segundo mais frequente em nível nacional ( 27,4 \%, correspondendo a 59 cursos ), e o conceito 4, o terceiro ( 20,5 \%, correspondendo a 44 cursos ).
Houve, ainda, 14 cursos ( 6,5 \% ) que receberam conceito 5, e 21 cursos ( 9,8 \% ) que receberam conceito 1.
Dos 215 cursos de Medicina Veterinária, nenhum ficou Sem Conceito ( SC ).
A região Norte participou com 17 cursos ou 7,9 \% do total nacional.
Desses, seis cursos receberam o conceito 2, o que equivale a 35,3 \% do total regional, além de ser o conceito modal na região, como já comentado.
A três cursos ( 17,6 \% ) atribuiu -se conceito 1, a cinco cursos ( 29,4 \% ) o conceito 3 e a três cursos ( 17,6 \% ), conceito 4.
Nenhum curso ficou Sem Conceito ( SC ) e nenhum curso recebeu o conceito 5 nessa região.
A região Nordeste participou com 31 cursos ou 14,4 \% do total nacional.
Nessa região, 16 cursos ( 51,6 \% em termos regionais ) obtiveram conceito 3, o conceito modal para a região, como já comentado.
Os conceitos 2 e 4 foram atribuídos, respectivamente, a cinco e a oito cursos ( 16,1 \% e 25,8 \%, respectivamente ).
Ainda nessa região, dois cursos ( 6,5 \% ) foram avaliados com o conceito 1.
Nenhum curso ficou Sem Conceito ( SC ) ou recebeu conceito 5.
Dos 93 cursos participantes da região Sudeste ( 43,3 \% do total nacional ), o conceito 3 foi o conceito modal, como já anteriormente assinalado, tendo sido obtido por 32 cursos ( 34,4 \% ).
O conceito 1 foi atribuído a seis cursos ( 6,5 \% ), e o conceito 2, a 30 cursos ( 32,3 \% ).
Receberam os conceitos 4 e 5, respectivamente, 16 cursos ( 17,2 \% ) e nove cursos ( 9,7 \% ).
Nessa região, nenhum curso ficou Sem Conceito ( SC ).
Os 51 cursos da região Sul corresponderam a 23,7 \% do total nacional.
Como já indicado, houve predominância do conceito 3, atribuído a 37,3 \% dos cursos da região ( conceito modal ), ou seja, atribuído a 19 dos 51 cursos participantes na região Sul.
O conceito 1 foi atribuído a três cursos ( 5,9 \% ), o conceito 2 foi atribuído a 14 cursos ( 27,5 \% ), o conceito 4 a dez cursos ( 19,6 \% ) e o 5 a cinco cursos ( 9,8 \% ).
Nenhum curso ficou Sem Conceito ( SC ) na região Sul.
Já dos 23 cursos participantes na região Centro-Oeste ( 10,7 \% do total nacional ), os conceitos 1 e 4 foram atribuídos a sete cursos cada ( 30,4 \% em termos regionais, para cada ), o conceito modal, como já fora destacado.
Quatro cursos ( 17,4 \% ) obtiveram conceito 2, e outros cinco ( 21,7 \% ), conceito 3.
Nenhum curso ficou Sem Conceito ( SC ) ou recebeu conceito 5.
No Gráfico 5.1, é apresentada a distribuição cumulativa do Conceito Enade segundo Grande Região.
Em um gráfico de distribuição cumulativa, em cada valor de abscissa ( eixo x ), o valor apresentado no eixo das ordenadas ( eixo y ) é o valor acumulado da variável ( conceito Enade ).
Observando a linha vertical pontilhada, que ajuda a visualizar o percentual de conceitos acumulados até o valor 3, ou seja, conceitos 1, 2 e 3, pode -se notar que, por exemplo, na região Norte ( linha azul ) mais de 80 \% dos cursos com conceito 3 ou abaixo deste valor.
As regiões Sul ( linha rosa ) e Centro-Oeste ( linha laranja ), por outro lado, apresentam em torno de 70 \% dos cursos na mesma situação : com conceito 3 ou abaixo.
Destaca -se que em algumas regiões o conceito 5 não foi alcançado por nenhum curso, como ocorre nas regiões Norte ( linha azul ), Nordeste ( linha verde ) e Centro-Oeste ( linha laranja ), nas quais todos os cursos têm conceito menor ou igual a 4, e as respectivas poligonais cumulativas alcançam 100 \% no conceito 4.
Regiões com melhores conceitos.
A região Norte ( linha azul ) apresenta os cursos com conceitos mais baixos, ao passo que a região Sul ( linha rosa ) apresenta a distribuição com valores maiores.
Cursos em IES situadas nas regiões Sudeste ( linha vermelha ), Nordeste ( linha verde ) e Centro-Oeste ( linha laranja ) apresentam uma situação intermediária. 5.2 CONCEITOS POR CATEGORIA ADMINISTRATIVA, POR MODALIDADE DE ENSINO E GRANDE REGIÃO Na Tabela 5.2, é apresentada a distribuição dos cursos de Medicina Veterinária participantes do Enade/2019, por Categoria Administrativa e Modalidade de ensino, de acordo com os conceitos por eles alcançados, segundo as Grandes Regiões brasileiras.
Dos 215 cursos participantes, 68 ( 31,6 \% ) eram ministrados em Instituições Públicas, e 147 ( 68,4 \% ), em Instituições Privadas.
Quanto à Modalidade de ensino, todos os cursos eram presenciais.
De acordo com as informações da Tabela 5.2, em termos nacionais, entre Instituições Públicas e Privadas, observa -se que, dos 14 cursos avaliados com conceito 5, todos eram oferecidos em IES Públicas.
Dos 68 cursos participantes de IES Públicas, o conceito 4 foi o valor modal, atribuído a 28 cursos ( 41,2 \% ).
Dos demais cursos de Instituições Públicas participantes, três cursos foram avaliados com conceito 1, quatro cursos foram avaliados com conceito 2, 19 cursos ( 27,9 \% ) foram avaliados com conceito 3 e, como já comentado, 14 cursos ( 20,6 \% ), com conceito 5.
Nenhum curso ficou Sem Conceito ( SC ). conceito 2, quatro cursos foram avaliados com conceito 3 e, como já comentado, 14 cursos ( 20,6 \% ), com conceito 5.
Nenhum curso ficou Sem Conceito ( SC ).
Na rede privada, o conceito modal foi 3, atribuído a 58 cursos, o correspondente a 39,5 \% dos 147 cursos da categoria.
Dos demais cursos participantes, 18 ( 12,2 \% ) receberam conceito 1, e 55 ( 37,4 \% ), conceito 2.
O conceito 4 foi atribuído a 16 cursos ( 10,9 \% ).
Nesta Categoria Administrativa, nenhum curso ficou Sem Conceito ( SC ) ou recebeu conceito 5.
Dos 215 cursos participantes da modalidade presencial, o conceito 3 foi o valor modal, atribuído a 77 cursos ( 35,8 \% ).
Dos demais cursos nesta modalidade de ensino, 21 cursos ( 9,8 \% ) receberam conceito 1, 59 cursos ( 27,4 \% ) receberam conceito 2, 44 cursos ( 20,5 \% ), conceito 4, e 14 cursos ( 6,5 \% ), conceito 5.
Nesta modalidade de ensino, nenhum curso ficou Sem Conceito ( SC ).
Na análise por região, observa -se que, na região Norte, as Instituições Privadas participaram com dez cursos ( 58,8 \% do total regional ), nenhum ficou Sem Conceito ( SC ), e nenhum recebeu os conceitos 4 e 5.
O conceito modal para as Instituições Privadas na região foi 2, atribuído a seis cursos, correspondendo a 60,0 \% dessa categoria na região.
Três cursos foram avaliados com conceito 1, e um curso, com conceito 3.
As Instituições Públicas participaram com sete cursos na região Norte ( 41,2 \% do total regional ), aos quais foram Na região Norte, as Instituições Privadas participaram com dez cursos ( 58,8 \% do total regional ), nenhum ficou Sem Conceito ( SC ), e nenhum recebeu os conceitos 4 e 5.
O conceito modal para as Instituições Privadas na região foi 2, atribuído a seis cursos, correspondendo a 60,0 \% dessa categoria na região.
Três cursos foram avaliados com conceito 1, e um curso, com conceito 3.
As Instituições Públicas participaram com sete cursos na região Norte ( 41,2 \% do total regional ), aos quais foram atribuídos os conceitos 3 e 4.
Nenhum curso ficou Sem Conceito ( SC ) ou com conceito 1, 2 e 5 nessa combinação de categoria e região.
Todos os cursos oferecidos na região Norte eram presenciais.
Na região Nordeste, a rede privada concentrou 12 dos 31 cursos participantes, o equivalente a 38,7 \% do total da região.
Dentre os cursos oferecidos por IES Privadas no Nordeste, cinco cursos foram avaliados com conceito 3, o conceito modal.
Dois cursos obtiveram conceito 1, três cursos receberam conceito 2 e dois cursos ficaram com conceito 4.
Nessa combinação de categoria e região, nenhum curso ficou Sem Conceito ( SC ) e nenhum recebeu conceito 5.
As Instituições Públicas da região Nordeste participaram com 19 cursos ( 61,3 \% do total da região ).
Desses, 11 obtiveram conceito 3, o valor modal, seguidos por outros seis cursos, que receberam conceito 4 e dois cursos com conceito 2.
Nenhum curso ficou Sem Conceito ( SC ) ou com os demais conceitos nessa região.
Todos os cursos oferecidos na região Nordeste eram presenciais.
Na região Sudeste, a proporção de cursos da rede privada, 81,7 \%, foi mais \textbf{elevada} do que a registrada nas demais regiões brasileiras, correspondendo a 76 dos 93 cursos participantes.
Nessa categoria e região, o conceito modal foi 3, atribuído a 32.
Os demais foram avaliados com conceito 1 ( cinco cursos ), conceito 2 ( 28 cursos ) e conceito 4 ( 11 cursos ).
Nenhum curso ficou Sem Conceito ( SC ) e nenhum recebeu conceito 5.
Dos 17 cursos oferecidos por Instituições Públicas na região Sudeste, o conceito modal foi 5, atribuído a nove cursos.
Nessa região, nenhum curso de IES Públicas ficou Sem Conceito ( SC ), um curso recebeu conceito 1, dois cursos receberam conceito 2 e cinco cursos receberam conceito 4.
Todos os cursos oferecidos na região Sudeste eram presenciais.
As Instituições Privadas concentraram 35 dos 51 cursos participantes da região Sul, 68,6 \% do total regional.
Desses, 17 cursos obtiveram conceito 3, o conceito modal.
Nesta combinação de Categoria Administrativa e Grande Região, três cursos receberam conceito 1, 14 cursos receberam conceito 2 e um curso recebeu conceito 4.
Nenhum curso recebeu conceito 5, e nenhum curso ficou Sem Conceito ( SC ).
As Instituições Públicas na região Sul participaram com 16 cursos ( 31,4 \% dos cursos da região ), aos quais foram atribuídos os conceitos 3, 4 e 5.
Nenhum curso da região Sul oferecido por IES Pública ficou Sem Conceito ( SC ) ou recebeu os conceitos 1 e 2.
Todos os cursos oferecidos na região Sul eram presenciais.
Na região Centro-Oeste, 14 dos 23 cursos participantes eram de Instituições Privadas ( 60,9 \% em termos regionais ).
Desses, cinco receberam conceito 1, o conceito modal.
Dos demais cursos, quatro receberam conceito 2, três receberam conceito 3 e dois receberam conceito 4.
Nenhum curso ficou Sem Conceito ( SC ), ou recebeu conceito 5.
Quanto aos nove cursos oferecidos por Instituições Públicas na região Centro-Oeste ( 39,1 \% do total regional ), cinco cursos apresentaram conceito 4, conceito modal, e os conceitos 1 e 3 foram atribuídos a dois cursos, cada.
Nenhum curso ficou Sem Conceito ( SC ) ou recebeu os demais conceitos nessa combinação de categoria e região.
Todos os cursos oferecidos na região Centro-Oeste eram presenciais.
No Gráfico 5.2, apresenta -se a distribuição do Conceito Enade, segundo a Categoria Administrativa da IES.
Os cursos em IES Públicas ( linha azul ) apresentam uma poligonal mais à direita do que os em IES Privadas ( linha verde ) e, consequentemente, uma distribuição de Conceitos Enade com valores maiores. 5.3 CONCEITOS POR ORGANIZAÇÃO ACADÊMICA E GRANDE REGIÃO Na Tabela 5.3, encontra -se a distribuição dos conceitos atribuídos aos cursos participantes do Enade/2019, na Área de Medicina Veterinária, por Organização Acadêmica, segundo as Grandes Regiões brasileiras.
Dos 215 cursos de Medicina Veterinária participantes, 113 eram oferecidos em Universidades, 59 em Centros Universitários, 36 em Faculdades e 7 em CEFET/IF.
Esta distribuição corresponde a, respectivamente, 52,6 \%, 27,4 \%, 16,7 \% e 3,3 \% dos cursos.
De acordo com os dados apresentados, de todos os 14 cursos avaliados com conceito 5, 13 eram vinculados a Universidades.
Os cursos vinculados a esse tipo de Organização Acadêmica tiveram o conceito 3 como conceito modal, atribuído a 46 cursos ( 40,7 \% ).
Os demais cursos vinculados a Universidades receberam os conceitos 1 ( três cursos ), 2 ( 19 cursos ), 4 ( 32 cursos ) e conceito 5 ( 13 cursos, como já mencionado ).
Nenhum curso ficou Sem Conceito ( SC ).
Dos cursos em Centros Universitários, o conceito modal foi 2, atribuído a 24 cursos ( 40,7 \% ).
Nenhum curso vinculado a esse tipo de Organização Acadêmica ficou Sem Conceito ( SC ), e os outros cursos receberam os conceitos 1 ( nove cursos ), 3 ( 20 cursos ) e 4 ( 6 cursos ).
Dos 36 cursos mantidos por Faculdades, 14 ( 38,9 \% ) receberam conceito 2, o conceito modal.
O segundo conceito mais frequentemente atribuído aos cursos oferecidos por essa categoria de IES foi o conceito 3, recebido por dez cursos.
Na sequência, nove cursos obtiveram conceito 1 e três cursos, conceito 4.
Dos cursos oferecidos em Faculdades, nenhum ficou Sem Conceito ( SC ) e nenhum recebeu conceito 5.
Dos sete cursos mantidos por CEFET/IF, três ( 42,9 \% ) receberam conceito 4, o conceito modal.
O segundo conceito mais frequentemente atribuído aos cursos oferecidos por essa categoria de IES foi o conceito 2, recebido por dois cursos.
Na sequência, um curso recebeu conceito 3 e um curso recebeu conceito 5.
Nenhum curso ficou Sem Conceito ( SC ), e nenhum curso recebeu conceito 1.
Considerando -se, separadamente, as regiões brasileiras, verifica -se que, na região Norte, as Universidades participaram com oito dos 17 cursos avaliados, ao qual foi atribuído conceito 3 a quatro cursos.
O conceito 4 foi atribuído a dois cursos, e os conceitos 1 e 2 foram atribuídos a um curso, cada.
Nenhum curso ficou Sem Conceito ( SC ) ou recebeu os demais conceitos nesta combinação de categoria e região.
Os Centros Universitários da região Norte foram representados por cinco cursos, o conceito 2 foi atribuído a três cursos, classe modal.
Os conceitos 1 e 3 foram atribuídos a um curso, cada.
Nenhum curso ficou Sem Conceito ( SC ) ou recebeu os demais conceitos nesta combinação de categoria e região.
As Faculdades participaram com três cursos na região Norte, os quais receberam os conceitos 1 ( um curso ) e 2 ( dois cursos, conceito modal ).
Nenhum curso ficou Sem Conceito ( SC ) ou recebeu os demais conceitos nesta combinação de categoria e região.
Já os CEFET/IF participaram com um curso na região Norte, ao qual foi atribuído o conceito 4.
Na região Nordeste, as Universidades participaram com 20 dos 31 cursos da Área de Medicina Veterinária oferecidos.
Nenhum curso ficou Sem Conceito ( SC ).
O conceito modal 3 foi atribuído a 13 cursos.
Seis cursos obtiveram conceito 4 e um, conceito 2.
Nenhum curso recebeu os demais conceitos nesta combinação de categoria e região.
Os Centros Universitários contaram com seis cursos participantes na região Nordeste, os conceitos 2 e 3 foram atribuídos a dois cursos cada, sendo as classes modais.
Para os demais cursos, foram atribuídos os conceitos 1 ( um curso ) e 4 ( um curso ).
Nenhum curso oferecido por Centros Universitários no Nordeste ficou Sem Conceito ( SC ) ou recebeu o conceito 5.
As Faculdades foram representadas por quatro cursos na região Nordeste, tendo sido atribuídos os conceitos 1, 2,3 e 4, a um curso, cada.
Na região Nordeste, nenhum curso mantido por Faculdades ficou Sem Conceito ( SC ) ou recebeu conceito 5.
Os CEFET/IF participaram com um curso na região Nordeste, avaliado com conceito 2.
Nenhum curso ficou Sem Conceito ( SC ) ou recebeu os demais conceitos nesta combinação de categoria e região.
Na região Sudeste, as Universidades concentraram 44 dos 93 cursos de Medicina Veterinária da região.
Dos cursos oferecidos em Universidades, nessa região, o conceito modal foi 3, atribuído a 15 cursos.
Os demais cursos receberam os conceitos 1 ( um curso ), 2 ( nove cursos ), 4 ( 11 cursos ) e 5 ( oito cursos ).
Nenhum curso ficou Sem Conceito ( SC ).
Os Centros Universitários participaram com 27 cursos na região Sudeste, dos quais 11 obtiveram o conceito modal 2.
O conceito 1 foi atribuído a dois cursos, o conceito 3, a dez cursos, e o conceito 4, a quatro cursos.
Nenhum curso ficou Sem Conceito ( SC ) ou recebeu conceito 5 nessa combinação de organização acadêmica e região.
Já as Faculdades foram representadas por 20 cursos na região Sudeste.
O conceito modal 2 foi atribuído a nove cursos.
Aos demais cursos foram atribuídos os conceitos 1 ( três cursos ), 3 ( sete cursos ) e 4 ( um curso ).
Nenhum curso ficou Sem Conceito ( SC ) ou recebeu conceito 5.
Os CEFET/IF participaram com dois cursos, nenhum ficando Sem Conceito ( SC ) e os dois cursos receberam os conceitos 2 e 5.
Dos 51 cursos da região Sul, 31 eram mantidos por Universidades, dos quais o conceito modal foi 3, atribuído a 12 cursos.
Os demais cursos receberam os conceitos 2 e 4 ( sete cursos cada ) e o conceito 5, cinco cursos.
Nesse tipo de organização, nenhum curso da região Sul ficou Sem Conceito ( SC ) ou recebeu conceito 1.
Dos 15 cursos participantes de Centros Universitários da região Sul, nenhum ficou Sem Conceito ( SC ), e nenhum recebeu o conceito 5.
Os conceitos 2 e 3, foram atribuídos a seis cursos, cada.
O conceito 1 foi atribuído a dois cursos e o conceito 4 atribuído a um curso.
Quanto aos três cursos vinculados a Faculdades na região Sul, os conceitos 1, 2 e 3 foram atribuídos a um curso, cada.
Nenhum curso ficou Sem Conceito ( SC ) ou recebeu conceito 4 ou 5.
Já os CEFET/IF participaram com dois cursos na região Norte, aos quais foi atribuído o conceito 4.
Na região Centro-Oeste, dez dos 23 cursos oferecidos eram mantidos por Universidades.
Para este tipo de Organização Acadêmica, nenhum curso ficou Sem Conceito ( SC ), os conceitos 1 e 2 foram atribuídos a um curso cada, o conceito 3 foi atribuído a dois cursos, e o conceito 4, conceito modal, foi atribuído a seis cursos.
Nesse tipo de organização, nenhum curso recebeu conceito 5.
Os Centros Universitários da região Centro-Oeste contaram com seis cursos.
O conceito 1 foi atribuído a três cursos, classe modal.
Os demais cursos receberam, conceito 2 ( dois cursos ) e conceito 3 ( um curso ).
Nenhum curso ficou Sem Conceito ( SC ) ou recebeu os demais conceitos nesta combinação de categoria e região.
Dos seis cursos oferecidos em Faculdades na região Centro-Oeste, três receberam o conceito modal 1.
Os conceitos 2, 3 e 4 foram atribuídos a um curso, cada.
Nesta região, nenhum curso oferecido em Faculdades ficou Sem Conceito ( SC ), e nenhum obteve conceito 5.
Já os CEFET/IF participaram com um curso na região Centro-Oeste, ao qual foi atribuído o conceito 3.
No Gráfico 5.3, apresenta -se a distribuição do Conceito Enade dos cursos de Medicina Veterinária, segundo a Organização Acadêmica da IES.
Os cursos em CEFET/IF ( linha laranja ) e Universidades ( linha azul ) apresentam uma poligonal mais à direita do que as demais, denotando uma melhor distribuição de Conceitos Enade.
No outro extremo, com distribuição muito semelhante entre si, encontram -se os cursos mantidos por Centros Universitários ( linha verde ) e Faculdades ( linha vermelha ), apresentando as respectivas poligonais mais à esquerda e, consequentemente, distribuição de Conceitos Enade com menores valores.

% matched lemmas: elevado, enfatizar
\end{textsample}
